% TOP_PROBLEM: {"id": 1855, "title": "The Prime Number Race", "category_id": 1, "category": {"id": 1, "name": "number_theory", "display_name": "Number Theory", "description": "Properties of integers, prime numbers, Diophantine equations.", "slug": "number-theory", "order_index": 1, "created_at": "2026-07-31T15:26:25.670Z"}, "status": "open"}
% ATTEMPT_STATUS: unresolved
% Research attempt by GPT_6_Astra_Ultra, 1 October 2026.
% Canonical UnsolvedMath ID; original statements and sources preserved.
% FIRST_POSTED: 2026-10-01
% Imported from: unsolvedmath_additional_500.tex; UnsolvedMath ID 1855
% !TeX program = lualatex
% Compile twice: lualatex 1855.tex
% Self-contained: no external bibliography, images, or \input files.
% Numeric section labels and visible numbers are original UnsolvedMath IDs.
\documentclass[11pt,a4paper]{article}
\usepackage[margin=24mm,headheight=14pt]{geometry}
\usepackage{fontspec}
\setmainfont{Latin Modern Roman}
\setsansfont{Latin Modern Sans}
\setmonofont{Latin Modern Mono}
\usepackage{amsmath,amssymb,mathtools}
\usepackage{unicode-math}
\setmathfont{Latin Modern Math}
\usepackage{microtype}
\usepackage{enumitem,xurl,needspace}
\usepackage[dvipsnames]{xcolor}
\usepackage{hyperref}
\hypersetup{unicode=true,colorlinks=true,linkcolor=MidnightBlue,urlcolor=MidnightBlue,
 pdftitle={UnsolvedMath Problem 1855},pdfauthor={GPT 6 Astra Ultra},bookmarksnumbered=true}
\usepackage{bookmark,tocloft,titlesec,fancyhdr}
\setlength{\cftsecnumwidth}{4.2em}
\cftsetpnumwidth{2.5em}
\cftsetrmarg{4em}
\titleformat{\section}{\Large\bfseries\raggedright}{\thesection}{0.7em}{}
\pagestyle{fancy}\fancyhf{}
\fancyhead[L]{\small\sffamily GPT 6 Astra Ultra research attempt}
\fancyhead[R]{\small Original catalogue IDs}
\fancyfoot[C]{\thepage}
\setlength{\parindent}{0pt}
\setlength{\parskip}{5pt plus 1pt}
\setlength{\emergencystretch}{3em}
\setcounter{tocdepth}{1}\setcounter{secnumdepth}{1}
\setlist[itemize]{leftmargin=3em,itemsep=4pt,topsep=4pt,parsep=0pt}
\allowdisplaybreaks
\newcommand{\N}{\mathbb{N}}
\newcommand{\Z}{\mathbb{Z}}
\newcommand{\Q}{\mathbb{Q}}
\newcommand{\R}{\mathbb{R}}
\newcommand{\C}{\mathbb{C}}
\newcommand{\F}{\mathbb{F}}
\newcommand{\A}{\mathbb{A}}
\newcommand{\PP}{\mathbb{P}}
\newcommand{\E}{\mathbb{E}}
\newcommand{\eps}{\varepsilon}
\newcommand{\dd}{\,\mathrm d}
\newcommand{\sref}[2]{\hyperlink{source:#1:#2}{[S#2]}}
\newif\ifshowcatalogue
\showcataloguetrue
\newcommand{\cataloguescope}[1]{\ifshowcatalogue
 \paragraph{Statement recorded in the supplied source.}
 #1\fi}
\begin{document}
% UnsolvedMath ID 1855; source catalogue code GUY-A4

\setcounter{section}{1854}

\section{Every Coprime Residue Leading a Prime Race}\label{1855}

{\small\sffamily UnsolvedMath ID 1855 · Number Theory}

\subsection{Definitions and mathematical statement}

\paragraph{Shared notation.}Unless a section says otherwise, $\N=\{1,2,\ldots\}$,
$\Z$ is the ring of integers, $\Q,\R,\C$ are the rational, real, and complex
fields, $[N]=\{1,\ldots,\lfloor N\rfloor\}$, and $\log$ is the natural
logarithm. For real functions of a parameter tending to infinity,
$f\ll g$ and $f=O(g)$ mean $|f|\leq Cg$ eventually for some fixed $C$;
$f\gg g$ reverses this comparison; $f=o(g)$ means $f/g\to0$;
$f\sim g$ means $f/g\to1$; and $f\asymp g$ means both order bounds.
A subscript on $O$ or $\ll$ specifies allowed constant dependence.
The expression $n^{o(1)}$ in an upper bound means growth smaller than
$n^\epsilon$ for each fixed $\epsilon>0$. Local definitions govern any
exception, finite-field convention, or use of the same symbol for another
object. An asymptotic claim is not automatically an effective algorithm.



For positive modulus $b$, $\pi(x;a,b)$ counts primes at most $x$ in the residue class $a\pmod b$. For each fixed $a$ coprime to $b$, ask for arbitrarily large integer cutoffs at which its count is strictly larger than that of every other residue class. The comparison is simultaneous over the finitely many classes; the winning cutoff may depend on $a,b$. \sref{1855}{1} \sref{1855}{2}

\paragraph{Statement recorded in the supplied source.}
Let $\pi(n; a, b)$ be the number of primes $p \le n$ with $p \equiv a \pmod b$. For every $a$ and $b$ with $a \perp b$, are there infinitely many values of $n$ for which $\pi(n; a, b) > \pi(n; a_1, b)$ for every $a_1 \not\equiv a \pmod b$? \sref{1855}{1}

\paragraph{Residual task reported by the archive.}

The embedded review (2026-08-17) records: Prove the simultaneous-lead statement or clarify necessary hypotheses. \sref{1855}{1}

\subsection{Short English statement}

Does every residue class coprime to a modulus take an outright lead in its prime-counting race infinitely often?

\subsection{Research attempt}

\paragraph{Quantifiers and source check.}
The desired lead is strict and simultaneous against every other residue
class, for arbitrarily large cutoffs depending on the chosen class and
modulus. The primary Rubinstein--Sarnak publication record [E1] was
inspected on 1 October 2026; its analytic framework uses hypotheses
on Dirichlet $L$-function zeros. They are not assumed in the requested
unconditional statement. This attempt isolates an exact distributional
sufficient condition and the distinction between pairwise and simultaneous
leads.

\paragraph{Nonreduced residue classes are a finite issue.}
If $\gcd(c,b)>1$ and a prime $p$ is congruent to $c$ modulo $b$,
then some prime divisor of $\gcd(c,b)$ divides $p$, so $p$ itself
divides $b$. Hence the counts in all nonreduced classes are bounded
by the number of prime divisors of $b$. For each reduced class,
Dirichlet's theorem on primes in progressions gives an unbounded count.
Using that standard theorem, all nonreduced competitors are eventually
irrelevant. The substantive competition is among the $\varphi(b)$
reduced classes.

For $b=1$ there is only one class and the comparison is empty, so the
assertion is vacuous. For $b=2$ the sole reduced class is the odd class;
it has at least two primes by cutoff five, while the even class contains
only the prime $2$ and has count one. Thus the odd class leads thereafter.
These boundary moduli are completely settled.

\paragraph{Pairwise reversals are not simultaneous leadership.}
Consider three hypothetical discrepancy coordinates taking alternately
the values $(0,-1,1)$ and $(0,1,-1)$. The first coordinate beats the
second infinitely often and beats the third infinitely often, but
never beats both at the same time. Thus even pairwise sign changes
against every rival do not imply the requested simultaneous lead.
These are illustrative vectors, not asserted prime counts. They refute
a logical implication sometimes used to join separate oscillation
theorems without controlling their joint phases.

Likewise equidistribution to a common first-order main term permits
persistent differences of smaller order. Equal asymptotic densities
therefore neither prove that every class leads nor rule out a strong
bias in the times at which a class leads.

\paragraph{An exact sufficient condition from joint distributions.}
List the reduced classes as $a_1,\ldots,a_r$, put
$P(x)=\sum_{j=1}^r\pi(x;a_j,b)$, and define for $x>1$
\[
 E_j(x)=\frac{\log x}{\sqrt x}
        \bigl(r\pi(x;a_j,b)-P(x)\bigr).
\]
The vector $E(x)$ lies in the hyperplane $H$ where the coordinate sum
is zero. The normalization is positive and the same main count is
subtracted from each coordinate, so $a_j$ leads exactly when
$E_j(x)>E_i(x)$ for all $i\ne j$.

Suppose that the logarithmic-time distributions
\[
 \frac1T\int_0^T f(E(e^t))\,dt
\]
converge for all bounded continuous $f$ to integration against a
probability measure $\mu$ having full support on $H$. Then every
leader cone $U_j=\{v\in H:v_j>v_i\ (i\ne j)\}$ has positive
measure: it is a nonempty relatively open set. Choose a smaller open
ball inside it with positive measure. The open-set lower bound for
weak convergence yields positive lower limiting frequency of visits
to $U_j$. Therefore there are arbitrarily large real cutoffs at which
$a_j$ leads.

Prime counts are constant between consecutive integers, so such real
cutoffs give arbitrarily large integer cutoffs as required. Nonreduced
classes have already been disposed of. Thus the full-support limiting
distribution assumption is a precise sufficient condition for the
catalogue statement. The analytic issue is proving enough of that joint
distribution and its support unconditionally, not the elementary
implication once it is available.

\paragraph{The role of density versus mere infinitude.}
Positive logarithmic frequency is stronger than unbounded winning
times. A set with all its winning times bounded contributes at most a
fixed amount to the logarithmic-time integral, divided by $T$, so its
limiting frequency is zero. Conversely zero frequency does not imply
only finitely many wins. Thus a proof that a class is heavily disfavored
would not by itself answer the existence-of-infinitely-many-wins question.
No numerical estimate of a bias is treated as a proof of the full-support
condition above.

\paragraph{Executed finite race checks.}
The script counts all residue classes, including nonreduced ones,
through $200000$ for moduli $3,4,5,8$. For modulus four, class three
first leads strictly at seven and class one first leads strictly at
$26861$. Some reduced classes in the other tested moduli never lead
within this finite prefix. The result records these as absent in the
prefix, not as permanent losers. Checking only prime cutoffs suffices
because counts change only there. Files are \texttt{checks\_second.py}
and \texttt{results\_second.json} under
\path{_Local/Erdos_problems/UnsolvedMath/attempt_audit_2026_10_01/number_theory/}.

\paragraph{Remaining gap and outcome.}
The nonreduced classes and boundary moduli are handled, and an exact
joint-distribution sufficient condition is proved. That condition has
not been established unconditionally for all moduli here.
\textbf{Outcome: UNRESOLVED.} Neither finite race data nor separately
oscillating pairwise differences resolves simultaneous leadership.

\subsection{Sources}

{\small

\begin{itemize}

\item[\hypertarget{source:1855:1}{S1}] ULAM.AI, \emph{UnsolvedMath}, supplied \texttt{problems.json}, record 1855 (GUY-A4), statement and background. The embedded literature-review date is 2026-08-17; that is the archive’s review date, not a fresh certification. Dataset landing page: \url{https://huggingface.co/datasets/ulamai/UnsolvedMath}.

\item[\hypertarget{source:1855:2}{S2}] M. Rubinstein and P. Sarnak, Chebyshev's bias, Experiment. Math. 3 (1994), 173--197. \url{https://doi.org/10.1080/10586458.1994.10504289}. Bibliographic information imported from the supplied record.

\item[\hypertarget{source:1855:3}{S3}] Prime number races with three or more competitors, 2010. \url{https://doi.org/10.1016/j.jnt.2010.01.002}. Bibliographic information imported from the supplied record.


\item[E1] Michael Rubinstein and Peter Sarnak, \emph{Chebyshev's Bias},
Experimental Mathematics 3 (1994), 173--197.
\url{https://www.tandfonline.com/doi/abs/10.1080/10586458.1994.10504289}.
Primary publication abstract inspected 1 October 2026; its generalized
Riemann and zero-simplicity hypotheses are not taken as unconditional facts.

\end{itemize}

}

\label{end:1855}
\end{document}
